\documentclass[11pt,a4paper]{article}

\usepackage[margin=1in]{geometry}
\usepackage{amsmath,amssymb}
\usepackage{booktabs}
\usepackage{graphicx}
\usepackage{siunitx}
\usepackage{multirow}
\usepackage{makecell}
\usepackage{hyperref}
\usepackage{natbib}
\usepackage{microtype}
\usepackage{xcolor}
\usepackage{caption}
\usepackage{subcaption}

\hypersetup{
  colorlinks=true,
  linkcolor=black,
  citecolor=blue!50!black,
  urlcolor=blue!50!black
}

\graphicspath{{figures/}}

\title{Downscaling Municipal Soybean Productivity Labels to Pixel Predictions with a Spectral--Temporal Transformer}
\author{}
\date{}

\begin{document}
\maketitle

\begin{abstract}
Official soybean productivity statistics in Brazil are reported at the municipality level, whereas Sentinel-2 observations resolve crop condition within municipalities.
We formulate this scale mismatch as a downscaling task: a spectral--temporal Transformer assigns a productivity contribution to each soybean pixel and is trained so that the municipal mean of those contributions matches IBGE PAM labels.
Using Sentinel-2 time series for Paran\'{a} (harvest years 2020--2024), augmented with gridded climate covariates, the trained pixel outputs provide intramunicipal productivity maps without field-level supervision.
On a 2021 spatial holdout, the climate-augmented model attains validation and test $R^2$ of $0.227$ and $0.228$.
Pooled early-season skill improves with additional months of imagery, peaking at $R^2=0.429$ after five season months, while the 2022 drought year remains difficult.
The results suggest that municipal labels can supervise spatially resolved productivity estimates, subject to remaining limits under extreme seasons and truncated inputs.
\end{abstract}

\section{Introduction}
\label{sec:introduction}

Soybean productivity maps at fine spatial resolution would support planning, insurance, and risk assessment, but field-level yield labels are rarely available at scale in Brazil.
What is available instead are official municipal statistics from the Brazilian Institute of Geography and Statistics (IBGE) Municipal Agricultural Production survey (PAM), together with dense Sentinel-2 time series that resolve spectral variation \emph{within} municipalities.
This paper treats that gap as a \emph{downscaling} task: learn pixel-level productivity contributions from municipal labels alone, then read out both municipal aggregates and intramunicipal maps.

Deep Transformers have been shown to extract informative representations from Sentinel-2 time series for crop mapping~\citep{xu2024sitsmoco}.
We adapt the Spectral--Temporal Network (STNet) from that line of research to continuous productivity regression (tonnes per hectare).
Each soybean pixel is scored by the network; municipal predictions are the mean of pixel scores and are supervised by PAM labels, using Sentinel-2 reflectance and Xavier climate covariates over Paran\'{a} for harvest years 2020--2024.

Our contributions are a municipal-to-pixel downscaling scheme that trains per-pixel productivity predictors under mean aggregation to IBGE PAM labels; a spectral--temporal Transformer for Paran\'{a} soybean that combines reflectance with Xavier climate covariates, with ablations that remove climate inputs or change temporal pooling; and a held-out evaluation on the 2021 growing season, including early-season truncations and qualitative intramunicipal maps.

\section{Data}
\label{sec:data}

Downscaling requires fine-resolution predictors and coarse labels.
Features are drawn from several sources (Table~\ref{tab:data-sources}): Sentinel-2 reflectance, Xavier climate covariates, and MapBiomas Solo soil properties at each soybean pixel, a MapBiomas soybean mask, and IBGE PAM productivity labels at the municipality level.

\begin{table}[t]
\centering
\caption{Summary of data sources used in the experiments.}
\label{tab:data-sources}
\footnotesize
\setlength{\tabcolsep}{3.5pt}
\begin{tabular}{@{}llp{2.35cm}ccp{2.35cm}@{}}
\toprule
Category & Variables & Unit & Related properties & Spatial res. & Source \\
\midrule
Satellite imagery
  & \makecell[l]{B2, B3, B4, B5, B6,\\B7, B8, B8A, B11, B12}
  & ---
  & Vegetation reflectance
  & 30\,m
  & Sentinel-2 L2A \\
\midrule
\multirow{6}{*}{Climate}
  & Cumulative precipitation
  & mm
  & Water availability
  & \multirow{6}{*}{$0.1^\circ$}
  & \multirow{6}{*}{\makecell[c]{Xavier\\(BR-DWGD)}} \\
\cmidrule(lr){2-4}
  & Max.\ dry-day streak
  & days
  & Drought stress
  &
  & \\
\cmidrule(lr){2-4}
  & Cumulative ETo
  & mm
  & Atmospheric demand
  &
  & \\
\cmidrule(lr){2-4}
  & Cumulative Rs
  & MJ\,m$^{-2}$
  & Solar radiation
  &
  & \\
\cmidrule(lr){2-4}
  & Cumulative Tmax
  & $^\circ$C$\cdot$day
  & Heat accumulation
  &
  & \\
\cmidrule(lr){2-4}
  & Cumulative Tmin
  & $^\circ$C$\cdot$day
  & Cold accumulation
  &
  & \\
\midrule
\multirow{4}{*}{Soil}
  & Clay (0--30\,cm)
  & \%
  & Texture
  & \multirow{4}{*}{30\,m}
  & \multirow{4}{*}{\makecell[c]{MapBiomas Solo\\Col.\ 3}} \\
\cmidrule(lr){2-4}
  & Silt (0--30\,cm)
  & \%
  & Texture
  &
  & \\
\cmidrule(lr){2-4}
  & Sand (0--30\,cm)
  & \%
  & Texture
  &
  & \\
\cmidrule(lr){2-4}
  & SOC stock (0--30\,cm)
  & \si{t/ha}
  & Organic carbon
  &
  & \\
\midrule
Land cover
  & Soybean mask (class 39)
  & ---
  & Crop mask
  & 30\,m
  & MapBiomas Col.\ 10 \\
\midrule
Labels
  & Soybean productivity
  & \si{t/ha}
  & Municipal supervision
  & Municipality
  & IBGE PAM \\
\bottomrule
\end{tabular}
\end{table}

\subsection{Study area and labels}
The study covers soybean-producing municipalities in the state of Paran\'{a}, Brazil.
Supervision uses municipal soybean productivity in tonnes per hectare (\si{t/ha}) from IBGE PAM, with municipality codes, harvest year, and a predefined validation and test partition within the holdout year.
No field- or pixel-level yield labels are used.
We retain municipality--year pairs whose Sentinel-2 coverage ratio relative to the reported planted area lies in $[0.8, 1.2]$, so that satellite coverage is roughly consistent with official area statistics.

\subsection{Imagery and climate covariates}
Predictors are per-pixel Sentinel-2 spectral time series, organized by municipality and growing season.
Each pixel sequence comprises ten spectral bands (B2, B3, B4, B5, B6, B7, B8, B8A, B11, B12) and a day-of-year index used for positional encoding.
Climate covariates come from the Xavier daily weather grid~\citep{xavier2022} and are temporally aligned to each Sentinel-2 observation date: cumulative precipitation from 1~October through the observation date; the maximum consecutive dry-day streak between successive Sentinel-2 dates (dry day: precipitation $<1$\,mm); and season-to-date cumulative reference evapotranspiration (ETo), solar radiation (Rs), maximum temperature (Tmax), and minimum temperature (Tmin).
All climate channels are sampled at the Xavier cell containing each pixel centroid and scaled before concatenation with normalized reflectance.
The proposed model uses all six Xavier climate channels together with the ten spectral bands (sixteen features per timestep).
Sequences are truncated or padded to a fixed length of $T=45$ timesteps.
Soybean pixels are selected with the MapBiomas Collection~10 land-cover product (class~39), which defines the intramunicipal support of the downscaled maps.
An ablation omits all climate covariates and uses the ten spectral bands alone.

\subsection{Train / validation / test split}
Splits follow a year-holdout design on municipalities.
All municipality--year observations outside the holdout year are used for training; municipalities in the holdout year are partitioned into validation and test sets.
The holdout year is \textbf{2021}.
After restricting to harvest years $\{2020,2021,2022,2023,2024\}$, the training set contains 997 municipality--years, while validation and test contain 132 and 123 municipalities in 2021, respectively.
Training target statistics used for loss normalization are mean $3.17$\,\si{t/ha} and standard deviation $0.93$\,\si{t/ha}.
Quantitative evaluation uses municipal aggregates of the pixel predictions; intramunicipal structure is examined qualitatively through pixel maps.

\section{Methods}
\label{sec:methods}

\subsection{Downscaling formulation}
Let a municipality $m$ in harvest year $y$ contain $N_m$ valid soybean pixels with time series $\mathbf{x}_p$.
A shared model $f_\theta$ maps each pixel to a scalar productivity contribution $\hat{y}_p = f_\theta(\mathbf{x}_p)$.
These pixel scores are the downscaled product of interest: after training they form intramunicipal productivity maps.
Municipal productivity, which is the only quantity for which labels exist, is recovered by mean aggregation,
\begin{equation}
  \hat{y}_{m,y} = \frac{1}{N_m}\sum_{p=1}^{N_m} \hat{y}_p ,
  \label{eq:mean-agg}
\end{equation}
and is supervised against the official PAM label $y_{m,y}$.
Thus learning proceeds from coarse labels to fine predictions: gradients update $f_\theta$ through the municipal mean, while inference retains the full pixel field $\{\hat{y}_p\}$.

\subsection{Architecture}
We use a regression adaptation of STNet~\citep{xu2024sitsmoco} as the pixel encoder.
Given pixel sequences with sixteen channels (ten spectral bands and six Xavier climate covariates), padding masks, and day-of-year indices, the network projects each timestep with a shared multilayer perceptron from 16 inputs through widths $32 \rightarrow 64 \rightarrow d_{\mathrm{model}}$, adds sinusoidal positional encodings indexed by day of year, applies a Transformer encoder with padding masks to model temporal dependencies, collapses the sequence to a single vector per pixel by temporal pooling, and maps the pooled embedding to a continuous scalar with a multilayer-perceptron head.
Temporal pooling uses NDVI-weighted averaging: temporal weights derived from NDVI are normalized and used as a soft collapse over time.
An alternative that averages only over valid (unmasked) timesteps, without NDVI weights, is evaluated as an ablation.
The spectral-only ablation uses the same architecture with ten-band inputs.

The proposed model uses one Transformer layer, $d_{\mathrm{model}}=128$, four attention heads, feed-forward width $d_{\mathrm{inner}}=128$, and dropout $0.3$.

\subsection{Loss and optimization}
Training minimizes mean squared error between aggregated municipal predictions and PAM targets after standardization by the training-set mean and standard deviation.
Reported metrics (RMSE, MAE, $R^2$, MAPE) are computed in original \si{t/ha} units on the municipal aggregates, which are the evaluable layer of the downscaling pipeline.

Optimization uses AdamW with learning rate $1\times10^{-4}$, weight decay $5\times10^{-4}$, and batches of 40 municipalities.
The learning rate is reduced at epochs 15 and 75; training stops early after ten epochs without improvement in validation $R^2$.
Because municipalities can contain many pixels, gradients within each municipal batch are accumulated over subsets of pixels.
No self-supervised pretraining is used: the network is trained from scratch under municipal supervision.

\section{Experiments}
\label{sec:experiments}

\subsection{Proposed model}
The proposed downscaling model combines Sentinel-2 reflectance with the full Xavier climate layout (precipitation plus cumulative ETo, Rs, Tmax, and Tmin) and uses one Transformer layer with $d_{\mathrm{model}}=128$, four attention heads, $d_{\mathrm{inner}}=128$, and dropout $0.3$.
Model selection maximizes validation $R^2$ of municipal aggregates; the best validation epoch is 12.
After selection, pixel scores from the same checkpoint are retained for intramunicipal maps.

\subsection{Architecture search without climate}
Separately, we grid-searched Transformer capacity and regularization using spectral inputs only:
$d_{\mathrm{model}}\in\{64,96,128\}$, number of heads $\in\{4,8\}$, $d_{\mathrm{inner}}\in\{64,128\}$, and dropout $\in\{0.3,0.4\}$, with a single Transformer layer.
Twenty-four configurations completed successfully.
The best spectral-only model ($d_{\mathrm{model}}=96$, eight heads, $d_{\mathrm{inner}}=64$, dropout $0.3$) is used in the spectral-only ablation.

\subsection{Evaluation protocol}
Held-out evaluation scores municipal means of pixel predictions against PAM labels on the 2021 validation and test municipalities.
We also report learning curves, predicted-versus-actual municipal scatters for harvest years 2020--2024, municipal error maps, and example intramunicipal pixel maps.

\subsection{Ablations}
We compare two ablations against the proposed model: a spectral-only setting that uses reflectance alone with the best spectral-only architecture from the grid search, and a temporal-pooling comparison that, with that spectral-only architecture fixed, contrasts NDVI-weighted pooling against mean pooling over valid timesteps.

\subsection{Early-season evaluation}
To assess whether downscaled estimates can be produced before the end of the season, we re-evaluate the proposed model using only the first $k$ months of imagery, $k=1,\ldots,6$.
Season months are counted from 1~October (month 1 $=$ October, \ldots, month 6 $=$ March).
For each $k$, observations in the retained window are truncated or padded to length $T=45$.
Metrics are reported on municipal aggregates, pooled over municipalities and years, and separately for 2021.

\section{Results}
\label{sec:results}

\subsection{Municipal aggregates of downscaled predictions}
Because pixel-level yield labels are unavailable, quantitative skill is assessed on municipal aggregates of the pixel predictions (Eq.~\ref{eq:mean-agg}).
Table~\ref{tab:main-metrics} summarizes validation and test metrics for the proposed spectral--climate downscaling model.
Validation $R^2$ peaks at epoch 12 ($0.227$), with RMSE $0.362$\,\si{t/ha} and MAPE $8.82\%$.
The corresponding test metrics are $R^2=0.228$, RMSE $0.351$\,\si{t/ha}, and MAPE $8.42\%$.
Absolute and relative municipal errors are therefore on the order of a few tenths of a tonne per hectare and about $8\%$--$9\%$ MAPE, with closely matched validation and test $R^2$ on the 2021 spatial holdout.

\begin{table}[t]
\centering
\caption{Validation and test metrics for municipal aggregates of the proposed downscaling model (best validation epoch 12). Units for RMSE/MAE are \si{t/ha}; MAPE is in percent.}
\label{tab:main-metrics}
\begin{tabular}{lcccc}
\toprule
Split & RMSE & MAE & $R^2$ & MAPE (\%) \\
\midrule
Validation (2021) & 0.362 & 0.287 & 0.227 & 8.82 \\
Test (2021)       & 0.351 & 0.271 & 0.228 & 8.42 \\
\bottomrule
\end{tabular}
\end{table}

Figure~\ref{fig:learning} shows learning dynamics.
Validation RMSE and $R^2$ improve through about epoch 12; after the scheduled learning-rate reduction near epoch 15, validation performance becomes unstable and eventually degrades, consistent with early stopping at epoch 22.

\begin{figure}[t]
\centering
\includegraphics[width=\textwidth]{learning_curves.pdf}
\caption{Training dynamics for the proposed spectral--climate downscaling model. The dashed line marks the best validation epoch (12).}
\label{fig:learning}
\end{figure}

\subsection{Architecture search without climate}
Table~\ref{tab:arch-top} lists the five best spectral-only configurations by validation $R^2$.
The top configuration is clearly ahead of the next-best models and is used as the spectral-only ablation below.
Wider spectral configurations often underperform, suggesting that capacity control and dropout matter for this dataset size and holdout design.

\begin{table}[t]
\centering
\caption{Top five spectral-only configurations by validation $R^2$ (24 completed runs). The proposed model instead uses the full Xavier climate layout with $d_{\mathrm{model}}=128$, four heads, $d_{\mathrm{inner}}=128$, and dropout $0.3$.}
\label{tab:arch-top}
\begin{tabular}{lcccccc}
\toprule
Rank & $d_{\mathrm{model}}$ & Heads & $d_{\mathrm{inner}}$ & Dropout & Val $R^2$ & Test $R^2$ \\
\midrule
1 & 96 & 8 & 64 & 0.3 & 0.165 & 0.117 \\
2 & 96 & 4 & 128 & 0.4 & 0.084 & 0.134 \\
3 & 64 & 4 & 64 & 0.4 & 0.076 & --- \\
4 & 64 & 4 & 128 & 0.3 & 0.071 & 0.119 \\
5 & 64 & 4 & 128 & 0.4 & 0.046 & 0.106 \\
\bottomrule
\end{tabular}
\end{table}

\subsection{Ablations}
\label{sec:ablations}

\paragraph{Spectral only.}
Training with spectral bands alone yields validation $R^2=0.165$ (RMSE $0.376$\,\si{t/ha}, MAPE $9.05\%$) and test $R^2=0.117$ (RMSE $0.375$\,\si{t/ha}, MAPE $9.09\%$).
Relative to the proposed spectral--climate model, removing climate covariates lowers both validation $R^2$ ($0.227\rightarrow 0.165$) and test $R^2$ ($0.228\rightarrow 0.117$), while MAPE remains near $9\%$.
Climate inputs therefore improve the municipal aggregates used for model selection and holdout evaluation under the 2021 spatial split.

\paragraph{Temporal pooling.}
With the best spectral-only architecture fixed, replacing NDVI-weighted temporal pooling by mean pooling over valid timesteps reduces validation $R^2$ from $0.165$ to $0.049$ (best epoch 12), while test $R^2$ remains nearly unchanged ($0.117$).
NDVI-weighted pooling therefore improves validation fit under the selection criterion used here, but does not translate into a clear test gain on this single holdout year.

\subsection{Per-year municipal diagnostics}
Figure~\ref{fig:scatter} shows predicted versus actual \emph{municipal} productivity for each harvest year under the proposed model.
Years 2020 and 2021 lie nearest the identity line (MAPE $8.1\%$--$8.6\%$), with 2021 $R^2=0.229$ consistent with the holdout evaluation.
Year 2022 remains difficult: RMSE rises to $1.14$\,\si{t/ha} and MAPE to $88.6\%$ despite near-zero $R^2$, reflecting a drought year with strong productivity collapse.
Years 2023--2024 are intermediate (Table~\ref{tab:per-year}).

\begin{figure}[t]
\centering
\includegraphics[width=\textwidth]{scatter_panel.pdf}
\caption{Predicted versus actual municipal soybean productivity (\si{t/ha}) by harvest year for the proposed downscaling model.
Points are colored by split (train / validation / test). The dashed line is identity.}
\label{fig:scatter}
\end{figure}

\begin{table}[t]
\centering
\caption{Per-year municipal metrics for aggregates of the proposed model (municipalities with coverage-filtered forecasts in each year).}
\label{tab:per-year}
\begin{tabular}{lccccc}
\toprule
Year & $n$ & RMSE & MAE & $R^2$ & MAPE (\%) \\
\midrule
2020 & 247 & 0.370 & 0.296 & 0.167 & 8.1 \\
2021 & 255 & 0.356 & 0.279 & 0.229 & 8.6 \\
2022 & 240 & 1.140 & 0.928 & $-0.001$ & 88.6 \\
2023 & 248 & 0.464 & 0.371 & 0.087 & 10.7 \\
2024 & 260 & 0.651 & 0.525 & 0.059 & 20.9 \\
\bottomrule
\end{tabular}
\end{table}

Figure~\ref{fig:map2021} maps municipal aggregate errors in the 2021 holdout year.
Figure~\ref{fig:guara} shows the corresponding intramunicipal product: pixel-level productivity contributions within Guarapuava, Paran\'{a}, obtained without field-level labels.
Spatial structure in such maps is the operational output of downscaling; municipal metrics only certify that the pixel field is consistent with PAM when averaged.

\begin{figure}[t]
\centering
\includegraphics[width=0.72\textwidth]{map_error_2021.png}
\caption{Municipal aggregate prediction error map for harvest year 2021 (holdout).}
\label{fig:map2021}
\end{figure}

\begin{figure}[t]
\centering
\includegraphics[width=0.72\textwidth]{guarapuava_heatmap.png}
\caption{Intramunicipal map of downscaled pixel-level productivity predictions for Guarapuava, Paran\'{a}.}
\label{fig:guara}
\end{figure}

\subsection{Early-season skill}
Table~\ref{tab:incomplete} and Figure~\ref{fig:incomplete} summarize early-season evaluation of municipal aggregates pooled over years.
With one or two season months, $R^2$ is negative and MAPE exceeds $34\%$.
Skill becomes clearly positive by month 3 ($R^2=0.266$) and peaks at month 5 ($R^2=0.429$, RMSE $0.648$\,\si{t/ha}, MAPE $26.2\%$), with a slight decline at month 6 ($R^2=0.411$).
Restricting the same protocol to 2021 alone, $R^2$ remains negative until month 6, where it reaches $0.229$ with RMSE $0.356$\,\si{t/ha}---matching the full-season holdout aggregate.
Early downscaled forecasts therefore improve as more of the season is observed when pooled across years, but the spatial holdout year requires nearly the full season window under this checkpoint.

\begin{table}[t]
\centering
\caption{Early-season evaluation of municipal aggregates (pooled years) as a function of season months $k$ retained from 1~October.}
\label{tab:incomplete}
\begin{tabular}{ccccc}
\toprule
$k$ & RMSE & MAE & $R^2$ & MAPE (\%) \\
\midrule
1 & 0.958 & 0.800 & $-0.247$ & 36.9 \\
2 & 0.907 & 0.793 & $-0.116$ & 34.6 \\
3 & 0.735 & 0.587 & 0.266 & 29.2 \\
4 & 0.680 & 0.523 & 0.373 & 27.3 \\
5 & 0.648 & 0.483 & 0.429 & 26.2 \\
6 & 0.658 & 0.477 & 0.411 & 26.9 \\
\bottomrule
\end{tabular}
\end{table}

\begin{figure}[t]
\centering
\includegraphics[width=\textwidth]{incomplete_series.pdf}
\caption{Early-season municipal metrics versus number of available season months $k$ (pooled years).}
\label{fig:incomplete}
\end{figure}

\section{Discussion}
\label{sec:discussion}

The central claim of this work is that municipal PAM labels are enough to train a pixel-level productivity model whose intramunicipal maps are the intended product, while municipal means of those maps remain evaluable against official statistics.
On the 2021 validation and test splits, those aggregates attain roughly $0.35$--$0.36$\,\si{t/ha} RMSE and about $8\%$--$9\%$ MAPE, with closely matched $R^2$ near $0.23$.
MAPE suggests that relative municipal errors are often practically small, whereas $R^2$ indicates that only a modest fraction of cross-municipality variance is explained on this single holdout year.
Pixel maps (e.g., Figure~\ref{fig:guara}) cannot be scored directly without field labels; their credibility rests on consistency of the municipal mean and on qualitative spatial structure.

Adding Xavier climate covariates improves both validation and test $R^2$ relative to the spectral-only ablation (validation $R^2$ $0.227$ vs.\ $0.165$; test $R^2$ $0.228$ vs.\ $0.117$).
Climate inputs are therefore useful for the municipal selection criterion and for spatial generalization within the holdout year under the protocol used here.

Architecture search on spectral inputs favored a compact Transformer ($d_{\mathrm{model}}=96$, modest feed-forward width, moderate dropout) for the spectral-only ablation, whereas the proposed model uses a wider configuration ($d_{\mathrm{model}}=128$, $d_{\mathrm{inner}}=128$, dropout $0.3$).
Combined with validation degradation after the learning-rate schedule step, this points to a regime where capacity, covariates, and optimization noise jointly affect year-holdout generalization of the downscaled aggregates.

Per-year diagnostics highlight 2022 as a severe failure mode: RMSE and MAPE rise sharply even though $R^2$ is near zero.
Interpreting $R^2$ alone can therefore be misleading under drought-driven level shifts; RMSE and MAPE remain essential, and models selected on a single holdout year may not transfer to extreme seasons---including the intramunicipal maps derived from the same checkpoint.

NDVI-weighted temporal pooling improved validation $R^2$ relative to mean pooling over valid timesteps on the spectral-only architecture, but test $R^2$ was essentially unchanged.
The ablation supports NDVI weighting as a useful inductive bias for model selection, while cautioning against overinterpreting a single holdout-year test difference.

Early-season results show progressive information gain as the season unfolds.
Pooled across years, municipal $R^2$ turns positive by month 3 and peaks near month 5 ($0.429$) before a slight decline at month 6.
On 2021 alone, truncated evaluations remain weak until the full six-month window recovers the holdout $R^2$.
Operational early warning with downscaled maps would therefore need models trained explicitly for truncated inputs, multi-year aggregation of forecasts, or stronger season-aware regularization.

A structural limitation is the mean-aggregation assumption: municipal productivity is treated as the average of pixel contributions, which ignores possible unequal area weighting, management heterogeneity, and label noise in PAM.
Without pixel-level references, we cannot quantify intramunicipal calibration error, only municipal consistency.
Further limitations include a single holdout year, the absence of external algorithmic baselines, and sensitivity to extreme seasons such as 2022.
Self-supervised pretraining was not evaluated here and remains future work.

\section{Conclusion}
\label{sec:conclusion}

We framed soybean productivity mapping in Paran\'{a} as downscaling: a spectral--temporal Transformer predicts productivity at every Sentinel-2 soybean pixel and is trained so that municipal means match IBGE PAM labels, yielding intramunicipal maps without field-level supervision.
On a 2021 spatial holdout, municipal aggregates of the climate-augmented model achieve approximately $8\%$--$9\%$ MAPE and $0.35$--$0.36$\,\si{t/ha} RMSE, with validation and test $R^2$ near $0.23$, and outperform a spectral-only ablation.
Early-season truncations improve pooled skill through about five season months (peak $R^2=0.429$), while the 2021 holdout needs nearly the full season window, and drought year 2022 remains challenging.
Future work includes multi-year holdouts, stronger baselines under identical splits, self-supervised pretraining, and pathways to validate intramunicipal structure when limited field references become available.


\section*{Data availability}
Municipal productivity labels are from IBGE PAM and provide the only supervised signal.
Sentinel-2 time series and Xavier climate covariates were processed to per-pixel sequences within municipalities, with coverage filtering as described in Section~\ref{sec:data}, enabling intramunicipal prediction maps after municipal training.
Further details are available from the authors upon reasonable request.

\bibliographystyle{plainnat}
\bibliography{refs}

\end{document}
